\documentclass[11pt,a4paper]{article}
\usepackage[margin=23mm]{geometry}
\usepackage{amsmath,amssymb,amsthm}
\setlength{\emergencystretch}{2em}
\newtheorem{theorem}{Theorem}
\newcommand{\gen}[1]{\langle#1\rangle}
\title{Completion of the critical six-point-head branch}
\author{Human Compression project: L044 continuation}
\date{}
\begin{document}\maketitle
Let $G\le\operatorname{Alt}(\Omega)$ be transitive with connected
string C-marking $\rho_0,\ldots,\rho_{r-1}$, all maximal omissions
intransitive, and $n=|\Omega|\ge12$. Assume the alternating rank
bound at every strictly smaller degree at least twelve. All complete
acquired proper primitive, anchored imprimitive and finite boundary
proofs remain retained. Suppose the first three actual marks form
an extendable supported six-point head from source Table~2(2) or (3).

\begin{theorem}
One has $r\le\lfloor(n-1)/2\rfloor$.
\end{theorem}
\begin{proof}
Suppose the bound fails. The complete critical-head reduction
gives $n=2r\ge18$, the literal supported head on $D$, its
actual natural alternating tail $G_{\ge3}$, and all head
and prefix interfaces used below. The full transitive-$B$
exclusion makes $B=G_{\ge4}$ intransitive on $\Omega\setminus D$.
Thus colour $3$ is paired; the exact early-cut calculation
also makes colour $4$ paired. The complete critical-tail
theorem supplies a unique ambient-eligible colour at least five.
Choose the least such colour $j$. Every colour $3,\ldots,j-1$
is therefore paired in the full ambient sense.

The pure-prefix unique-cut lemma gives the two actual $G_j$
orbits $C=C_{j-1}$ and $A=\Omega\setminus C$, where $C_k$
is the orbit of $P_k=G_{\le k}$ containing $D$. Every mark
after $j$ fixes $C$ pointwise. Any shared label lies before
$j$, and its first index $g$, if present, is at least three.

If the cut is unshared, its two groups are the actual supported
transitive prefix and suffix. The complete later-unshared theorem
leaves only an actual opposite six-point head, or $j=r-1$
with $A=\{z\}$ and $P_{r-2}=\operatorname{Alt}(\Omega\setminus\{z\})$.
The first case is excluded by the complete opposite-head theorem.
In the second, minimality gives paired colours $3,\ldots,r-2$.
The complete terminal-singleton reduction applies with these exact
hypotheses and produces an actual opposite head, again impossible.

It remains that a first shared label $g$ exists.
The complete interior and endpoint first-shared-head proofs
exclude $g=3$, using their respective, distinct support interfaces.
The complete first-shared-four proof excludes $g=4$.
Hence $g\ge5$.

Suppose $P_{g-1}$ is transitive on its support. Reverse the
entire string and apply the complete transitive-tail alternative
at the unique cut. In that reversed orientation the crossing
is $i=r-1-j$, its last shared label is $h=r-1-g<r-1$,
and its supported transitive tail is exactly $P_{g-1}$.
If $i>0$, the preceding group is the original $G_{>j}$,
which fixes $C$ pointwise by the unique-cut lemma. Thus the
chosen-orientation hypothesis is satisfied literally.
The exact index-location lemma says that
its first-head exit has $g=3$, already excluded, while its
other exit is an actual opposite head with $g=r-3$.
The opposite-head theorem excludes that alternative too.

We may therefore assume $P_{g-1}$ is intransitive on its
support. If $j=r-1$, the complete endpoint equality route
applies: $B$ is intransitive, $j$ is the least later unique
colour, and $g\ge5$. It proves $g=r-3$ and an actual pure
opposite Table~2(2) head, already excluded.

For $j<r-1$, the complete first-later-cut equality theorem
has all its hypotheses. It gives
\[
 j=g+1,\quad t=r-g-1\ge2,\quad |A|=2t,\quad
 \operatorname{Fix}(P_{g-1})=A\cup\{b\},
\]
where $b\in C$ is the crossing endpoint. The support of
$P_{g-1}$ has exactly two nontrivial orbits, the head orbit
and one exterior orbit, of total size $2g+1$.
The suffix $Q=G_{>j}$ has rank $t-1$ and is either
transitive on $A$ or has two faithful natural $S_t$ orbits.
The complete transitive-equality-suffix theorem excludes
the first possibility.

In the second possibility $u=\rho_g$ exchanges the two
$Q$-orbits equivariantly: the earlier prefix fixes $A$
pointwise and $\gen{u,Q}$ is transitive on that actual
omission orbit. The complete pair-path argument gives the
literal marking
\[
 q_i=(a_{i-1}\ a_i)(a'_{i-1}\ a'_i),\qquad 1\le i<t,
 \qquad u|_A=\prod_{i=0}^{t-1}(a_i\ a'_i).
\]
For precision, every pair-image omission is intransitive,
since adjoining the actual $u$ would otherwise make $A$
one orbit without its omitted, supported mark. One selected
edge per label is a tree on $t$ pairs; its omission cuts force
each image to be precisely its own edge transposition, and
string commutation makes the tree a path. The endpoint anchor
fixes its orientation. This argument includes $t=2,3$ and
does not import an intersection property for a quotient.

The complete small-equality-suffix lemma proves that the
crossing $s=\rho_{g+1}$ fixes the old prefix support pointwise,
so $s=(b\ a_0)w$ with $w$ an odd involution on
$A\setminus\{a_0\}$. If $t\ge4$, the complete diagonal
suffix theorem forces $s$ to fix that entire latter set;
then $s$ is a single transposition, impossible in $G$.
If $t=3$, the complete six-point suffix theorem makes the
last ambient omission transitive, also impossible.

For $t=2$, the complete four-point reduction leaves its
one literal decorated form. The complete quotient criterion
and its application exclude it: they give a genuine string
C-marking of rank $g+1$ on the natural alternating group
of degree $2g+2<n=2g+6$. Here $g\ge5$, so the explicit
smaller-degree induction applies and gives rank at most $g$.
The retained proof checks the complete quotient kernel, full IP
and every nonidentity projected mark; no quotient assertion is
inferred from an abstract group identification.

Every least later cut and every actual prefix/suffix orbit
case has now been exhausted. The contrary rank range is
impossible, proving the theorem.
\end{proof}

\paragraph{Attribution and boundary.}
This assembles the exceptional-head route underlying
Propositions~5.28--5.32 of Cameron--Fernandes--Leemans--Mixer,
\emph{Highest rank of a polytope for $A_n$}, arXiv:1605.09173v1.
The source cut strategy and all complete acquired replacement
proofs remain paid. The assembly uses the general equality route;
the additional separately proved exclusions of individual later
first-shared indices are unnecessary premises here.
This proves the critical head branch, with its induction explicit.
The full main-theorem assembly must still account for primitive
and imprimitive omissions, the all-intransitive branches without
an exceptional head, finite bases and the lower-bound construction.
No whole-original, question, numerical or formal credit is claimed.
\end{document}
